%% ─────────────────────────────────────────────────────────────────────────────
%% CV — Talvin Ackbaraly (Version française)
%% Serif (XCharter), one restrained accent colour, thin rules, bullets, no icons
%% Compile: latexmk -pdf cv-talvin-ackbaraly-fr.tex
%% ─────────────────────────────────────────────────────────────────────────────
\documentclass[10pt, a4paper]{extarticle}

\usepackage[T1]{fontenc}
\usepackage[utf8]{inputenc}
\usepackage[french]{babel}
\usepackage[top=1.2cm, bottom=1cm, left=1.8cm, right=1.8cm]{geometry}
\usepackage{XCharter}
\usepackage[letterspace=80]{microtype}
\usepackage{titlesec}
\usepackage{enumitem}
\usepackage{tabularx}
\usepackage{ragged2e}
\usepackage{xcolor}
\usepackage[scaled=0.94,noupquote]{inconsolata}
\usepackage{tikz}
\usepackage[hidelinks]{hyperref}

\hypersetup{pdftitle={CV — Talvin Ackbaraly}, pdfauthor={Talvin Ackbaraly}}

%% ── Colour system: one accent, everything else neutral ───────────────────────
\definecolor{accent}{HTML}{2563A0}  % THE accent: name, sections, links, bullets
\definecolor{ink}{HTML}{1A1A1A}     % body text (warm near-black, not pure 000)
\definecolor{soft}{HTML}{6E6E6E}    % secondary information (context, tech)
\definecolor{hair}{HTML}{C9C9C9}    % section rules
\definecolor{chipbg}{HTML}{EAF0F7}  % tag pill background
\definecolor{chipfg}{HTML}{2F5B87}  % tag pill text

\hypersetup{colorlinks=true, urlcolor=accent, linkcolor=accent}
\AtBeginDocument{\color{ink}}

\pagestyle{empty}
\setlength{\parindent}{0pt}
\setlength{\parskip}{0pt}
\linespread{1.02}
\RaggedRight

%% Section: letter-spaced small caps over a hairline
\titleformat{\section}
  {\large\scshape\lsstyle\color{accent}}{}{0pt}{}
  [\vspace{-5pt}{\color{hair}\rule{\textwidth}{0.4pt}}]
\titlespacing*{\section}{0pt}{12pt}{5pt}

\frenchsetup{StandardItemLabels=true}
\setlist[itemize]{leftmargin=1.1em, label={\color{accent!70}\textbullet}, topsep=3pt, itemsep=2.5pt, parsep=0pt}

%% Tag pill: \chip{CUDA}
\newcommand{\chip}[1]{%
  \tikz[baseline=(c.base)]{\node[inner xsep=4.5pt, inner ysep=2pt, rounded corners=2.5pt,
    fill=chipbg, text=chipfg, font=\footnotesize] (c) {#1};}\hspace{1pt}%
}
%% \entry{Title}{Date}{Subtitle}{Tags}
\newcommand{\entry}[4]{%
  \textbf{#1} \hfill #2\\
  \textit{#3} \hfill {#4}\par
}
%% \project{Name}{Context}{Tags}
\newcommand{\project}[3]{%
  \textbf{#1}{\color{soft}\enspace—\enspace\textit{#2}} \hfill {#3}\par
}
\newcommand{\gap}{\par\vspace{6pt}}

\begin{document}

%% ══ EN-TÊTE ══════════════════════════════════════════════════════════════════
\begin{center}
  {\color{accent}\fontsize{25}{30}\selectfont\scshape\textls[40]{Talvin Ackbaraly}}\\[6pt]
  Élève ingénieur — Vision par ordinateur \& calcul GPU\\[5pt]
  {\small
  \href{mailto:talvin.ackbaraly@icloud.com}{talvin.ackbaraly@icloud.com} \enspace·\enspace
  \href{tel:+33769389008}{07 69 38 90 08} \enspace·\enspace
  \href{https://talvin-ackbaraly.com}{talvin-ackbaraly.com}\\[2pt]
  \href{https://github.com/talvinckb}{github.com/talvinckb} \enspace·\enspace
  \href{https://www.linkedin.com/in/talvin-ackbaraly}{linkedin.com/in/talvin-ackbaraly}\\[5pt]
  \textit{Recherche un stage de fin d'études de 6 mois à partir de février 2027}}
\end{center}

%% ══ FORMATION ════════════════════════════════════════════════════════════════
\section{Formation}

\entry{EPITA — École d'ingénieurs en informatique}{sept. 2022 -- 2027}
      {Diplôme d'ingénieur, majeure Image}{\small\color{soft}\textit{Paris, France}}
{\small\color{soft}Traitement \& synthèse d'images, vision par ordinateur, deep learning, calcul GPU.}
\gap
\entry{Université du Québec à Chicoutimi (UQAC)}{janv. -- mai 2024}
      {Semestre d'échange académique}{\small\color{soft}\textit{Québec, Canada}}

%% ══ PROJETS ══════════════════════════════════════════════════════════════════
\section{Projets}

\project{Détection d'illustrations pour la BnF}{projet de fin d'études, équipe de 4, en cours}{\chip{PyTorch}\chip{YOLO26}\chip{Florence-2}}
\begin{itemize}
  \item Conception d'une chaîne de \textbf{détection et de classification d'illustrations} dans des documents numérisés.
  \item Entraînement et benchmark de modèles de détection et de classification (YOLO26, Florence-2) sur \textbf{30\,Go}.
  \item Définition des métriques d'évaluation avec le client et présentation des résultats en points hebdomadaires.
\end{itemize}
\gap
\project{Détection de mouvement temps réel sur GPU}{équipe de 4}{\chip{CUDA}\chip{C++}\chip{Nsight}}
\begin{itemize}
  \item Portage CPU $\rightarrow$ GPU d'une chaîne de détection de mouvement vidéo : \textbf{accélération $\times$24}.
  \item Profilage et suppression des goulots d'étranglement avec NVIDIA Nsight.
\end{itemize}
\gap
\project{Simulation de fluides temps réel}{équipe de 2}{\chip{C++}\chip{OpenGL}\chip{GLSL}}
\begin{itemize}
  \item Moteur SPH simulant \textbf{75\,000 particules à plus de 60 FPS}, avec calcul physique en compute shaders.
\end{itemize}
\gap
\project{Reconnaissance automatique de plaques d'immatriculation}{équipe de 2}{\chip{C++}\chip{Python}\chip{OpenCV}}
\begin{itemize}
  \item Chaîne complète (prétraitement, extraction de candidats, HOG, Random Forest) : \textbf{précision de 86\,\%}.
  \item Benchmark des implémentations Python et C++ ; gestion via GitLab (issues, MR, CI).
\end{itemize}
\gap
\project{Prédiction de la fonction pulmonaire à partir de scanners CT}{équipe de 4}{\chip{Python}\chip{XGBoost}\chip{React}\chip{Docker}}
\begin{itemize}
  \item Prédiction de la capacité vitale forcée par segmentation pulmonaire et ML classique (XGBoost, sans CNN).
  \item Intégration dans une application web (front-end React), conteneurisée avec Docker et un pipeline GitLab CI.
\end{itemize}

%% ══ EXPÉRIENCE ═══════════════════════════════════════════════════════════════
\section{Expérience professionnelle}

\entry{Studiocanal}{sept. 2025 -- janv. 2026}{Stage — Développeur Full-Stack \& DevOps}{\chip{Java}\chip{Spring Boot}\chip{Angular}\chip{AWS}}
\begin{itemize}
  \item Nouvelles fonctionnalités pour une plateforme interne (droits, contrats, catalogue) utilisée par \textbf{100+ collaborateurs}.
  \item Mise en place du pipeline CI/CD GitLab d'un nouveau module ; déploiements en DEV, UAT et PROD.
  \item Intervention sur l'infrastructure AWS (EC2, RDS, IAM) au sein d'une équipe Agile.
\end{itemize}
\gap
\entry{Hitachi Solutions France}{juin -- juil. 2022}{Stage — Consultant informatique}{\chip{Python}\chip{Django}}
\begin{itemize}
  \item Prototype d'application de gestion fournisseurs, produits et contrats pour Sonepar, en méthode Agile SAFe.
\end{itemize}

%% ══ COMPÉTENCES ══════════════════════════════════════════════════════════════
\section{Compétences}

{\small\renewcommand{\arraystretch}{1.12}
\begin{tabularx}{\textwidth}{@{}l@{\hspace{1.5em}}>{\RaggedRight\arraybackslash}X@{}}
  \textbf{\color{accent}Vision par ordinateur \& IA} & Traitement d'images, deep learning, PyTorch, OpenCV, YOLO, Florence-2, XGBoost \\
  \textbf{\color{accent}GPU \& graphisme}            & CUDA, Nsight, OpenGL, GLSL, compute shaders \\
  \textbf{\color{accent}Programmation}               & C, C++, Python, Java, TypeScript, SQL \\
  \textbf{\color{accent}Frameworks}                  & Spring Boot, Quarkus, FastAPI, Django, Angular, React \\
  \textbf{\color{accent}Outils}                      & Linux, Git, Docker, GitLab CI, AWS, PostgreSQL, MongoDB, SQLite \\
  \textbf{\color{accent}Langues}                     & Français (langue maternelle), anglais (professionnel, TOEIC 845)
\end{tabularx}}

\end{document}
